% Level 06 - Model lifecycle
% Question answered: how does a model get trained, judged, promoted and watched, and how does the system learn from analysts?
% Facts: report MLOps scope (DVC + MLflow), Beta promotion gate (F1 >= 0.80, FPR < 0.03 percent), review records as feedback, synthetic pretraining then real fine-tuning.
\documentclass[tikz,border=10pt]{standalone}
% =============================================================================
% MethaneWatch diagram style  (shared by every diagram under diagrams/)
% -----------------------------------------------------------------------------
% Why one shared file: every diagram must look like part of the same set, and a
% colour change (for example a layer re-coloured) should be one edit, not 33.
%
% Usage, at the top of any diagram:
%   \documentclass[tikz,border=10pt]{standalone}
%   % =============================================================================
% MethaneWatch diagram style  (shared by every diagram under diagrams/)
% -----------------------------------------------------------------------------
% Why one shared file: every diagram must look like part of the same set, and a
% colour change (for example a layer re-coloured) should be one edit, not 33.
%
% Usage, at the top of any diagram:
%   \documentclass[tikz,border=10pt]{standalone}
%   % =============================================================================
% MethaneWatch diagram style  (shared by every diagram under diagrams/)
% -----------------------------------------------------------------------------
% Why one shared file: every diagram must look like part of the same set, and a
% colour change (for example a layer re-coloured) should be one edit, not 33.
%
% Usage, at the top of any diagram:
%   \documentclass[tikz,border=10pt]{standalone}
%   \input{../_common/mwstyle.tex}
%   \begin{document} ... \end{document}
%
% Colour meaning is fixed across the whole set (see diagrams/README.md). Two colours plus grey:
%   blue = MethaneWatch's own parts (application, data, models, science, services)
%   red  = evidence, security, audit and hard limits (the things that must never be bypassed)
%   grey = people, cloud and infrastructure, external systems, notes
% The older names green, orange, purple, teal and gold are kept as aliases so that no diagram
% source needs to change; they now resolve to blue (green, orange, purple) or grey (teal, gold).
% =============================================================================
\usepackage[T1]{fontenc}
\usepackage[scaled=0.95]{helvet}
\renewcommand{\familydefault}{\sfdefault}
\hyphenpenalty=10000 \exhyphenpenalty=10000 % never hyphenate words inside diagrams
\usepackage{pifont}
\usepackage{amssymb}
\usepackage{array}
\usepackage{booktabs}
\usetikzlibrary{positioning,arrows.meta,shapes.geometric,shapes.multipart,%
  shapes.misc,fit,backgrounds,calc,decorations.pathreplacing,matrix,patterns}

% ---- palette ----------------------------------------------------------------
\definecolor{mwink}{HTML}{1F2933}
\definecolor{mwgrey}{HTML}{6B7785}
\definecolor{mwline}{HTML}{4A5560}
\definecolor{mwlight}{HTML}{F2F4F7}
\definecolor{mwblue}{HTML}{1F6FB2}   \definecolor{mwblueL}{HTML}{E3EFFA}
\definecolor{mwred}{HTML}{B5483F}    \definecolor{mwredL}{HTML}{F8E3E0}
% aliases (see note above)
\colorlet{mwgreen}{mwblue}   \colorlet{mwgreenL}{mwblueL}
\colorlet{mworange}{mwblue}  \colorlet{mworangeL}{mwblueL}
\colorlet{mwpurple}{mwblue}  \colorlet{mwpurpleL}{mwblueL}
\colorlet{mwteal}{mwgrey}    \colorlet{mwtealL}{mwlight}
\colorlet{mwgold}{mwgrey}    \colorlet{mwgoldL}{mwlight}

% ---- node and edge styles ---------------------------------------------------
\tikzset{
  >={Stealth[length=2.2mm]},
  every node/.style={font=\sffamily\small,text=mwink},
  box/.style={rectangle,rounded corners=2.5pt,draw=mwline,line width=0.7pt,
    fill=white,align=center,inner sep=5pt,minimum height=7mm},
  blue/.style={draw=mwblue,fill=mwblueL},
  green/.style={draw=mwgreen,fill=mwgreenL},
  orange/.style={draw=mworange,fill=mworangeL},
  purple/.style={draw=mwpurple,fill=mwpurpleL},
  red/.style={draw=mwred,fill=mwredL},
  teal/.style={draw=mwteal,fill=mwtealL},
  gold/.style={draw=mwgold,fill=mwgoldL},
  grey/.style={draw=mwgrey,fill=mwlight},
  strong/.style={line width=1.6pt},
  ext/.style={box,dashed,draw=mwgrey,fill=mwlight},
  actor/.style={box,gold,rounded corners=8pt,minimum height=9mm},
  store/.style={cylinder,shape border rotate=90,aspect=0.22,draw=mwgreen,fill=mwgreenL,
    line width=0.7pt,align=center,inner sep=3pt,minimum height=11mm,minimum width=20mm},
  group/.style={rectangle,rounded corners=4pt,draw=mwline,dashed,line width=0.7pt,
    inner sep=7pt},
  glabel/.style={anchor=north west,font=\sffamily\bfseries\footnotesize,text=mwgrey,inner sep=2pt},
  note/.style={font=\sffamily\footnotesize,text=mwgrey,align=left,inner sep=1pt},
  tiny/.style={font=\sffamily\scriptsize},
  head/.style={font=\sffamily\bfseries\small},
  flow/.style={->,line width=0.8pt,draw=mwline},
  dflow/.style={->,line width=0.8pt,draw=mwline,dashed},
  bflow/.style={<->,line width=0.8pt,draw=mwline},
  lab/.style={font=\sffamily\scriptsize,text=mwgrey,fill=white,inner sep=1pt},
}

% ---- title / footer helpers -------------------------------------------------
% \dtitle{Finding or name}{One-line subtitle: what the diagram answers}
\newcommand{\dtitle}[2]{%
  \node[anchor=south west,font=\sffamily\bfseries\Large,text=mwink,inner sep=0pt] (dt)
    at ([yshift=12mm]current bounding box.north west) {#1};
  \node[anchor=north west,font=\sffamily\footnotesize,text=mwgrey,inner sep=0pt]
    at ([yshift=-1.2mm]dt.south west) {#2};}
% \dfoot{Source or reading note}
\newcommand{\dfoot}[2][14cm]{%
  \node[anchor=north west,note,text width=#1] at
    ([yshift=-4mm]current bounding box.south west) {#2};}
% in-node text helpers: \nm{Name} bold name, \tech{technology} small grey italic
\newcommand{\nm}[1]{{\bfseries\small #1}}
\newcommand{\tech}[1]{{\scriptsize\itshape\color{mwgrey} #1}}
% check mark used in matrices
\newcommand{\yes}{\ding{51}}

%   \begin{document} ... \end{document}
%
% Colour meaning is fixed across the whole set (see diagrams/README.md). Two colours plus grey:
%   blue = MethaneWatch's own parts (application, data, models, science, services)
%   red  = evidence, security, audit and hard limits (the things that must never be bypassed)
%   grey = people, cloud and infrastructure, external systems, notes
% The older names green, orange, purple, teal and gold are kept as aliases so that no diagram
% source needs to change; they now resolve to blue (green, orange, purple) or grey (teal, gold).
% =============================================================================
\usepackage[T1]{fontenc}
\usepackage[scaled=0.95]{helvet}
\renewcommand{\familydefault}{\sfdefault}
\hyphenpenalty=10000 \exhyphenpenalty=10000 % never hyphenate words inside diagrams
\usepackage{pifont}
\usepackage{amssymb}
\usepackage{array}
\usepackage{booktabs}
\usetikzlibrary{positioning,arrows.meta,shapes.geometric,shapes.multipart,%
  shapes.misc,fit,backgrounds,calc,decorations.pathreplacing,matrix,patterns}

% ---- palette ----------------------------------------------------------------
\definecolor{mwink}{HTML}{1F2933}
\definecolor{mwgrey}{HTML}{6B7785}
\definecolor{mwline}{HTML}{4A5560}
\definecolor{mwlight}{HTML}{F2F4F7}
\definecolor{mwblue}{HTML}{1F6FB2}   \definecolor{mwblueL}{HTML}{E3EFFA}
\definecolor{mwred}{HTML}{B5483F}    \definecolor{mwredL}{HTML}{F8E3E0}
% aliases (see note above)
\colorlet{mwgreen}{mwblue}   \colorlet{mwgreenL}{mwblueL}
\colorlet{mworange}{mwblue}  \colorlet{mworangeL}{mwblueL}
\colorlet{mwpurple}{mwblue}  \colorlet{mwpurpleL}{mwblueL}
\colorlet{mwteal}{mwgrey}    \colorlet{mwtealL}{mwlight}
\colorlet{mwgold}{mwgrey}    \colorlet{mwgoldL}{mwlight}

% ---- node and edge styles ---------------------------------------------------
\tikzset{
  >={Stealth[length=2.2mm]},
  every node/.style={font=\sffamily\small,text=mwink},
  box/.style={rectangle,rounded corners=2.5pt,draw=mwline,line width=0.7pt,
    fill=white,align=center,inner sep=5pt,minimum height=7mm},
  blue/.style={draw=mwblue,fill=mwblueL},
  green/.style={draw=mwgreen,fill=mwgreenL},
  orange/.style={draw=mworange,fill=mworangeL},
  purple/.style={draw=mwpurple,fill=mwpurpleL},
  red/.style={draw=mwred,fill=mwredL},
  teal/.style={draw=mwteal,fill=mwtealL},
  gold/.style={draw=mwgold,fill=mwgoldL},
  grey/.style={draw=mwgrey,fill=mwlight},
  strong/.style={line width=1.6pt},
  ext/.style={box,dashed,draw=mwgrey,fill=mwlight},
  actor/.style={box,gold,rounded corners=8pt,minimum height=9mm},
  store/.style={cylinder,shape border rotate=90,aspect=0.22,draw=mwgreen,fill=mwgreenL,
    line width=0.7pt,align=center,inner sep=3pt,minimum height=11mm,minimum width=20mm},
  group/.style={rectangle,rounded corners=4pt,draw=mwline,dashed,line width=0.7pt,
    inner sep=7pt},
  glabel/.style={anchor=north west,font=\sffamily\bfseries\footnotesize,text=mwgrey,inner sep=2pt},
  note/.style={font=\sffamily\footnotesize,text=mwgrey,align=left,inner sep=1pt},
  tiny/.style={font=\sffamily\scriptsize},
  head/.style={font=\sffamily\bfseries\small},
  flow/.style={->,line width=0.8pt,draw=mwline},
  dflow/.style={->,line width=0.8pt,draw=mwline,dashed},
  bflow/.style={<->,line width=0.8pt,draw=mwline},
  lab/.style={font=\sffamily\scriptsize,text=mwgrey,fill=white,inner sep=1pt},
}

% ---- title / footer helpers -------------------------------------------------
% \dtitle{Finding or name}{One-line subtitle: what the diagram answers}
\newcommand{\dtitle}[2]{%
  \node[anchor=south west,font=\sffamily\bfseries\Large,text=mwink,inner sep=0pt] (dt)
    at ([yshift=12mm]current bounding box.north west) {#1};
  \node[anchor=north west,font=\sffamily\footnotesize,text=mwgrey,inner sep=0pt]
    at ([yshift=-1.2mm]dt.south west) {#2};}
% \dfoot{Source or reading note}
\newcommand{\dfoot}[2][14cm]{%
  \node[anchor=north west,note,text width=#1] at
    ([yshift=-4mm]current bounding box.south west) {#2};}
% in-node text helpers: \nm{Name} bold name, \tech{technology} small grey italic
\newcommand{\nm}[1]{{\bfseries\small #1}}
\newcommand{\tech}[1]{{\scriptsize\itshape\color{mwgrey} #1}}
% check mark used in matrices
\newcommand{\yes}{\ding{51}}

%   \begin{document} ... \end{document}
%
% Colour meaning is fixed across the whole set (see diagrams/README.md). Two colours plus grey:
%   blue = MethaneWatch's own parts (application, data, models, science, services)
%   red  = evidence, security, audit and hard limits (the things that must never be bypassed)
%   grey = people, cloud and infrastructure, external systems, notes
% The older names green, orange, purple, teal and gold are kept as aliases so that no diagram
% source needs to change; they now resolve to blue (green, orange, purple) or grey (teal, gold).
% =============================================================================
\usepackage[T1]{fontenc}
\usepackage[scaled=0.95]{helvet}
\renewcommand{\familydefault}{\sfdefault}
\hyphenpenalty=10000 \exhyphenpenalty=10000 % never hyphenate words inside diagrams
\usepackage{pifont}
\usepackage{amssymb}
\usepackage{array}
\usepackage{booktabs}
\usetikzlibrary{positioning,arrows.meta,shapes.geometric,shapes.multipart,%
  shapes.misc,fit,backgrounds,calc,decorations.pathreplacing,matrix,patterns}

% ---- palette ----------------------------------------------------------------
\definecolor{mwink}{HTML}{1F2933}
\definecolor{mwgrey}{HTML}{6B7785}
\definecolor{mwline}{HTML}{4A5560}
\definecolor{mwlight}{HTML}{F2F4F7}
\definecolor{mwblue}{HTML}{1F6FB2}   \definecolor{mwblueL}{HTML}{E3EFFA}
\definecolor{mwred}{HTML}{B5483F}    \definecolor{mwredL}{HTML}{F8E3E0}
% aliases (see note above)
\colorlet{mwgreen}{mwblue}   \colorlet{mwgreenL}{mwblueL}
\colorlet{mworange}{mwblue}  \colorlet{mworangeL}{mwblueL}
\colorlet{mwpurple}{mwblue}  \colorlet{mwpurpleL}{mwblueL}
\colorlet{mwteal}{mwgrey}    \colorlet{mwtealL}{mwlight}
\colorlet{mwgold}{mwgrey}    \colorlet{mwgoldL}{mwlight}

% ---- node and edge styles ---------------------------------------------------
\tikzset{
  >={Stealth[length=2.2mm]},
  every node/.style={font=\sffamily\small,text=mwink},
  box/.style={rectangle,rounded corners=2.5pt,draw=mwline,line width=0.7pt,
    fill=white,align=center,inner sep=5pt,minimum height=7mm},
  blue/.style={draw=mwblue,fill=mwblueL},
  green/.style={draw=mwgreen,fill=mwgreenL},
  orange/.style={draw=mworange,fill=mworangeL},
  purple/.style={draw=mwpurple,fill=mwpurpleL},
  red/.style={draw=mwred,fill=mwredL},
  teal/.style={draw=mwteal,fill=mwtealL},
  gold/.style={draw=mwgold,fill=mwgoldL},
  grey/.style={draw=mwgrey,fill=mwlight},
  strong/.style={line width=1.6pt},
  ext/.style={box,dashed,draw=mwgrey,fill=mwlight},
  actor/.style={box,gold,rounded corners=8pt,minimum height=9mm},
  store/.style={cylinder,shape border rotate=90,aspect=0.22,draw=mwgreen,fill=mwgreenL,
    line width=0.7pt,align=center,inner sep=3pt,minimum height=11mm,minimum width=20mm},
  group/.style={rectangle,rounded corners=4pt,draw=mwline,dashed,line width=0.7pt,
    inner sep=7pt},
  glabel/.style={anchor=north west,font=\sffamily\bfseries\footnotesize,text=mwgrey,inner sep=2pt},
  note/.style={font=\sffamily\footnotesize,text=mwgrey,align=left,inner sep=1pt},
  tiny/.style={font=\sffamily\scriptsize},
  head/.style={font=\sffamily\bfseries\small},
  flow/.style={->,line width=0.8pt,draw=mwline},
  dflow/.style={->,line width=0.8pt,draw=mwline,dashed},
  bflow/.style={<->,line width=0.8pt,draw=mwline},
  lab/.style={font=\sffamily\scriptsize,text=mwgrey,fill=white,inner sep=1pt},
}

% ---- title / footer helpers -------------------------------------------------
% \dtitle{Finding or name}{One-line subtitle: what the diagram answers}
\newcommand{\dtitle}[2]{%
  \node[anchor=south west,font=\sffamily\bfseries\Large,text=mwink,inner sep=0pt] (dt)
    at ([yshift=12mm]current bounding box.north west) {#1};
  \node[anchor=north west,font=\sffamily\footnotesize,text=mwgrey,inner sep=0pt]
    at ([yshift=-1.2mm]dt.south west) {#2};}
% \dfoot{Source or reading note}
\newcommand{\dfoot}[2][14cm]{%
  \node[anchor=north west,note,text width=#1] at
    ([yshift=-4mm]current bounding box.south west) {#2};}
% in-node text helpers: \nm{Name} bold name, \tech{technology} small grey italic
\newcommand{\nm}[1]{{\bfseries\small #1}}
\newcommand{\tech}[1]{{\scriptsize\itshape\color{mwgrey} #1}}
% check mark used in matrices
\newcommand{\yes}{\ding{51}}

\begin{document}
\begin{tikzpicture}[
  c/.style={box,text width=3.0cm,align=left,inner sep=4pt,minimum height=1.9cm,font=\sffamily\footnotesize}]

  % layer bands
  \node[note,anchor=east,text width=2.2cm,align=right] at (-0.9,3.2) {\bfseries Data layer};
  \node[note,anchor=east,text width=2.2cm,align=right] at (-0.9,0.0) {\bfseries Model layer};
  \node[note,anchor=east,text width=2.2cm,align=right] at (-0.9,-3.2) {\bfseries Serving layer};
  \draw[draw=mwline!40] (-3.2,4.5) -- (21.5,4.5) (-3.2,1.6) -- (21.5,1.6) (-3.2,-1.6) -- (21.5,-1.6) (-3.2,-4.8) -- (21.5,-4.8);

  % data layer
  \node[c,green] (dv) at (1.8,3.2) {\nm{Versioned data}\\ DVC: recipes, stacks, manifests with hashes};
  \node[c,green] (sp) at (6.0,3.2) {\nm{Fixed splits}\\ Grouped by event, site, date; real test sealed};
  \node[c,green,dashed] (hn) at (14.4,3.2) {\nm{Reviewed cases}\\ Rejected and Uncertain candidates become hard negatives};

  % model layer
  \node[c,orange] (tr) at (1.8,0) {\nm{Train}\\ Synthetic pretrain, then real fine-tune; config and seed logged};
  \node[c,orange] (ev) at (6.0,0) {\nm{Evaluate}\\ IoU, F1, false-positive rate, per-site, on real test};
  \node[c,orange] (rg) at (10.2,0) {\nm{Register}\\ MLflow run, data version, metrics, model card};
  \node[c,red,strong] (gt) at (14.4,0) {\nm{Release gate}\\ Beta: F1 $\ge$ 0.80 and FPR $<$ 0.03\%; else hold};

  % serving layer
  \node[c,blue] (dp) at (14.4,-3.2) {\nm{Promote}\\ Approved flag set; API loads that version};
  \node[c,blue] (sv) at (10.2,-3.2) {\nm{Serve}\\ Inference worker on the approved model};
  \node[c,purple] (mo) at (6.0,-3.2) {\nm{Monitor}\\ Candidate rate, review outcomes, drift per region};
  \node[c,teal,dashed] (rb) at (18.6,-3.2) {\nm{Roll back}\\ Previous approved version in 60 s};

  \draw[flow] (dv) -- (sp); \draw[flow] (dv.south) -- (tr.north); \draw[flow] (sp.south) -- (ev.north);
  \draw[flow] (tr) -- (ev); \draw[flow] (ev) -- (rg); \draw[flow] (rg) -- (gt);
  \draw[flow] (gt) -- node[lab,right] {pass} (dp);
  \draw[flow] (dp) -- (sv); \draw[flow] (sv) -- (mo);
  \draw[dflow] (dp) -- (rb); \draw[dflow] (rb.north) |- node[lab,pos=0.25,right] {failure or regression} (gt.east);

  % feedback loops
  \draw[flow,mwred,line width=1.1pt] (mo.west) -| (tr.south);
  \node[lab,rotate=90,text=mwred] at (1.5,-2.35) {loop 2};
  \draw[flow,mwpurple,line width=1.1pt] (mo.north) -- ++(0,0.8) -| node[lab,pos=0.25,above,text=mwpurple] {loop 1: analyst review} (hn.south west);
  \draw[flow,mwpurple,line width=1.1pt] (hn.west) -- (sp.east);

  \dtitle{Two loops keep the model honest: analysts feed it, monitoring triggers it}
         {Level 06 - lifecycle. Three layers, one gate. Loop 1 turns reviews into training cases; loop 2 retrains when monitoring shows drift or false alarms.}
  \dfoot[22cm]{Alpha targets (IoU 0.70, F1 0.75) are research results to measure and report; the Beta gate decides promotion. The two are different and never interchangeable. Rollback in 60 s is a Beta target.}
\end{tikzpicture}
\end{document}
